\begin{enumerate}
\item On calcule les premiers termes de la suite :
\[
u_1=2u_0+1=2\times 5+1=11,
\]
\[
u_2=2u_1+1=2\times 11+1=23,
\]
\[
u_3=2u_2+1=2\times 23+1=47.
\]
L'affirmation 1 est donc \textbf{vraie}.

\bigskip

\item Perdre un dixième de ses abonnés revient à conserver les neuf dixièmes, c'est-à-dire multiplier par :
\[
1-\dfrac{1}{10}=\dfrac{9}{10}=0{,}9.
\]
Puis Mathétix gagne 200 abonnés. On a donc, pour tout entier naturel $n$ :
\[
V_{n+1}=0{,}9V_n+200.
\]
La formule proposée $V_{n+1}=1{,}1V_n+200$ correspondrait à une augmentation de $10\,\%$, et non à une perte d'un dixième.

L'affirmation 2 est donc \textbf{fausse}.

\bigskip

\item Le premier terme de la suite $(w_n)$ est : $w_0=\dfrac{0+4}{3}=\dfrac{4}{3}$, et non 4.

L'affirmation 3 est donc \textbf{fausse}.

\bigskip

\item La probabilité qu'un élève choisi au hasard se rende au lycée en bus est :
\[P(B)=\dfrac{150}{200}=\dfrac{15}{20}=\dfrac{3}{4}.\]

L'affirmation 4 est donc \textbf{vraie}.

\bigskip

\item On trouve d'abord la valeur manquante dans case du tableau : il y a 150 élèves qui prennent le bus, dont 90 filles. Il y a donc $150-90=60$ garçons qui prennent le bus.

La probabilité cherchée est donc :
\[P_B(G)=\dfrac{60}{150}=\dfrac{2}{5}=0{,}4\neq0{,}6.\]

L'affirmation 5 est donc \textbf{fausse}.
\end{enumerate}
